% --- Archivo: 02_dualidad_general_wolfe.tex ---

% =========================================================
\section{Dualidad para un problema general}
% =========================================================

El análisis de dualidad se extiende a problemas con restricciones no lineales de igualdad y desigualdad mediante la función Lagrangiana. Dado el problema general de optimización, al cual llamaremos \textbf{problema primal}:
\begin{mini}{x}{f(x)}{\label{eq:primal_general}}{}
    \addConstraint{x \in D = \{x \in \Omega \mid h(x) = 0, \, g(x) \le 0\} }
\end{mini}
donde $\Omega \subset \Rn$, $f \colon \Omega \to \R$, $h \colon \Omega \to \R^l$ y $g \colon \Omega \to \R^m$.

\begin{definition}[Función Lagrangiana]\label{def:lagrangiana_general}
    Asociada al problema primal \eqref{eq:primal_general}, se define la \textbf{función Lagrangiana} con respecto a las restricciones funcionales como la aplicación $L \colon \Omega \times \R^l \times \R^m \to \R$ dada por:
    \begin{equation}
        L(x, \lambda, \mu) = f(x) + \interno{\lambda}{h(x)} + \interno{\mu}{g(x)},
    \end{equation}
    donde los vectores $\lambda \in \R^l$ y $\mu \in \R^m$ reciben el nombre de \textbf{multiplicadores de Lagrange}.
\end{definition}

% ---------------------------------------------------------
\subsection{Formulación Minimax del Problema Primal}
% ---------------------------------------------------------

La Lagrangiana penaliza implícitamente la violación de las restricciones funcionales al buscar el peor escenario respecto a las variables duales.

\begin{lemma}[Penalización implícita]\label{lem:sup_lagrangiana}
    Para cualquier $x \in \Omega$, se verifica:
    \begin{equation}
        \sup_{\substack{\lambda \in \R^l \\ \mu \ge 0}} L(x, \lambda, \mu) =
        \begin{cases}
            f(x)    & \text{si } x \in D, \\
            +\infty & \text{si } x \in \Omega \setminus D.
        \end{cases}
        \label{eq:sup_lagrangiana}
    \end{equation}
\end{lemma}

\begin{proof}
    Si $x \in D$, entonces $h(x) = 0$ y $g(x) \le 0$. Como $\mu \ge 0$, el término $\interno{\mu}{g(x)} \le 0$. El valor máximo de la Lagrangiana se alcanza de forma única con $\mu = 0$ (o en cualquier multiplicador si $g_i(x)=0$), resultando $\sup L = f(x)$.

    Si $x \notin D$, alguna restricción se infringe. Si $h_j(x) \neq 0$, basta con tomar $\lambda_j = t \cdot \sgn(h_j(x))$ con $t \to \infty$. Si $g_i(x) > 0$, tomamos $\mu_i = t \to \infty$. En ambos casos, el supremo diverge a $+\infty$.
\end{proof}

\begin{remark}
    El \cref{lem:sup_lagrangiana} permite reescribir de forma equivalente el problema primal \eqref{eq:primal_general} como un problema (tipo minimax) sin restricciones funcionales, sobre el conjunto base $\Omega$:
    \begin{equation}
        \min \left\{ \sup_{(\lambda, \mu) \in \R^l \times \R^m_+} L(x, \lambda, \mu) \right\} \quad \text{sujeto a} \quad x \in \Omega.
        \label{eq:minimax_primal_problem}
    \end{equation}
    Además, el \textbf{valor óptimo primal} $\bar{v}$ queda caracterizado por la ecuación:
    \begin{equation}
        \bar{v} = \inf_{x \in \Omega} \left\{ \sup_{(\lambda, \mu) \in \R^l \times \R^m_+} L(x, \lambda, \mu) \right\},
        \label{eq:minimax_primal_value}
    \end{equation}
    donde es obligatorio el uso del operador ínfimo ($\inf$), ya que, en ausencia de compacidad en $\Omega$, el mínimo global podría no alcanzarse.
\end{remark}

\begin{example}[Evaluación del supremo de la Lagrangiana]\label{ex:penalizacion_lagrangiana}
    Dado el problema en $\Omega = \R^2$:
    \begin{mini*}{ }{x_1 + x_2}{}{ }
        \addConstraint{x_1 ^2 + x_2 ^2 - 1}{= 0}
        \addConstraint{-x_1}{\leq  0}
    \end{mini*}
    la Lagrangiana asociada es $L(x, \lambda, \mu) = x_1 + x_2 + \lambda(x_1^2 + x_2^2 - 1) - \mu x_1$.
    \begin{enuthm}
        \item Para el punto factible $x^A = (1, 0)^\top \in D$, se tiene $L(x^A, \lambda, \mu) = 1 - \mu$. Su supremo sobre $\mu \ge 0$ se alcanza en $\bar{\mu} = 0$, resultando $f(x^A) = 1$.
        \item Para el punto infactible de igualdad $x^B = (2, 0)^\top$, se tiene $L(x^B, \lambda, \mu) = 2 + 3\lambda - 2\mu$. El supremo respecto a $\lambda \in \R$ diverge trivialmente a $+\infty$.
        \item Para el punto infactible de desigualdad $x^C = (-1, 0)^\top$, se tiene $L(x^C, \lambda, \mu) = -1 + \mu$. Como la variable $\mu$ puede ser arbitrariamente grande, el supremo sobre $\mu \ge 0$ es $+\infty$.
    \end{enuthm}
\end{example}

% ---------------------------------------------------------
\subsection{El Problema Dual General}
% ---------------------------------------------------------

El \textbf{problema dual} se define alternando el orden de los operadores de optimización en la formulación minimax del primal \eqref{eq:minimax_primal_problem}:
\begin{equation}
    \bar{w} = \max_{\substack{\lambda \in \R^l \\ \mu \ge 0}} \left\{ \inf_{x \in \Omega} L(x, \lambda, \mu) \right\}.
    \label{eq:minimax_dual}
\end{equation}

Para simplificar su análisis, introducimos la \textbf{función dual} $\varphi : \R^l \times \R^m \to [-\infty, +\infty)$ como el ínfimo parcial respecto a las variables primales:
\begin{equation}
    \varphi(\lambda, \mu) = \inf_{x \in \Omega} L(x, \lambda, \mu).
\end{equation}
El dominio de factibilidad dual es el conjunto $\Delta = \{(\lambda, \mu) \in \R^l \times \R^m \mid \mu \ge 0, \, \varphi(\lambda, \mu) > -\infty\}$, de modo que el problema dual se expresa como:
\begin{equation}
    \max \varphi(\lambda, \mu) \quad \text{sujeto a} \quad (\lambda, \mu) \in \Delta.
    \label{eq:dual_general}
\end{equation}

\begin{theorem}[Propiedades de la función dual]\label{thm:propiedades_dual}
    El conjunto de factibilidad dual $\Delta$ es convexo y la función dual $\varphi$ es cóncava sobre $\Delta$.
\end{theorem}

\begin{proof}
    Para cada punto fijo $x \in \Omega$, la función de multiplicadores $(\lambda, \mu) \mapsto L(x, \lambda, \mu)$ es afín y, por ende, cóncava. Como la función dual $\varphi$ es el ínfimo puntual de una familia de funciones afines continuas, conserva la concavidad de estas. La convexidad del dominio $\Delta$ se deduce de forma inmediata al ser el dominio de definición de una función cóncava propia.
\end{proof}

El problema dual siempre es un problema de optimización convexa (maximización de una función cóncava sobre un conjunto convexo), con independencia de las propiedades topológicas o de convexidad del problema primal original.

% ---------------------------------------------------------
\subsection{Teorema de Dualidad Débil}
% ---------------------------------------------------------

\begin{theorem}[Teorema de Dualidad Débil]\label{thm:dualidad_debil_general}
    Para cualquier punto factible primal $x \in D$ y cualquier punto factible dual $(\lambda, \mu) \in \Delta$, se cumple:
    \begin{equation}
        \varphi(\lambda, \mu) \le f(x).
    \end{equation}
\end{theorem}

\begin{proof}
    Evaluando el ínfimo parcial en el punto particular $x \in D$:
    \[
        \varphi(\lambda, \mu) = \inf_{z \in \Omega} L(z, \lambda, \mu) \le L(x, \lambda, \mu).
    \]
    Como $h(x) = 0$ y $\interno{\mu}{g(x)} \le 0$ para todo $\mu \ge 0$, se deduce que $L(x, \lambda, \mu) \le f(x)$.
\end{proof}

\begin{example}[Obtención analítica de la función dual]\label{ex:propiedades_dual}
    Sea el problema de optimización:
    \[
        \min \quad x^2 \quad \text{sujeto a} \quad x - 2 \le 0.
    \]
    La Lagrangiana para $\mu \ge 0$ es $L(x, \mu) = x^2 + \mu x - 2\mu$. Resolvemos el ínfimo sobre $\Omega = \R$ igualando la derivada respecto a $x$ a cero:
    \[
        2x + \mu = 0 \implies x^*(\mu) = -\frac{\mu}{2}.
    \]
    Sustituyendo $x^*$ en la Lagrangiana, se obtiene la función dual:
    \[
        \varphi(\mu) = -\frac{\mu^2}{4} - 2\mu.
    \]
    Dado que el ínfimo es finito para todo $\mu$, el dominio dual es $\Delta = [0, +\infty)$, el cual es convexo. Además, $\varphi''(\mu) = -1/2 < 0$, confirmando que la función dual es estrictamente cóncava sobre su dominio.
\end{example}

\begin{definition}[]
    La diferencia entre el valor óptimo primal $\bar{v}$ y el dual $\bar{w}$ se llama \textbf{salto de dualidad} ($\bar{v} - \bar{w} \ge 0$).
\end{definition}

\begin{corollary}[Ausencia de salto de dualidad]\label{cor:no_gap}
    Si existen puntos factibles $\bar{x} \in D$ y $(\bar{\lambda}, \bar{\mu}) \in \Delta$ tales que $f(\bar{x}) = \varphi(\bar{\lambda}, \bar{\mu})$, entonces $\bar{x}$ y $(\bar{\lambda}, \bar{\mu})$ son soluciones óptimas globales de sus respectivos problemas, y el salto de dualidad es nulo.
\end{corollary}

% ---------------------------------------------------------
\subsection{Patologías y Casos Límite en la Teoría de Dualidad}
% ---------------------------------------------------------

Los siguientes tres ejemplos ilustran cómo la falta de convexidad o la estructura de los dominios pueden impedir el alcance de los óptimos o dar lugar a saltos de dualidad estrictamente positivos.

\begin{example}[Mínimo primal no alcanzado, sin salto de dualidad]\label{ex:no_primal_sol}
    Sea el problema de optimización definido sobre el conjunto base abierto $\Omega = (0, +\infty)$:
    \begin{mini*}{ }{1 / x}{}{ }
        \addConstraint{-x}{\le 0}
    \end{mini*}
    El dominio factible es $D = (0, +\infty)$. El valor del ínfimo primal es $\bar{v} = 0$, pero no se alcanza para ningún $x$ finito. El problema primal no tiene solución.
    
    Para cualquier $\mu \ge 0$, la Lagrangiana es $L(x, \mu) = 1/x - \mu x$. Evaluamos la función dual $\varphi(\mu) = \inf_{x > 0} (1/x - \mu x)$:
    \begin{itemize}
        \item Si $\mu > 0$, tomando $x \to \infty$ el término $-\mu x$ decae sin límite inferior, resultando $\varphi(\mu) = -\infty$.
        \item Si $\mu = 0$, se tiene $\varphi(0) = \inf_{x > 0} 1/x = 0$.
    \end{itemize}
    Así, el dominio dual es el conjunto unitario $\Delta = \{0\}$. El problema dual consiste en evaluar $\varphi(0) = 0$, con valor óptimo dual $\bar{w} = 0$. No existe salto de dualidad ($\bar{v} = \bar{w} = 0$) y el dual admite solución óptima en $\bar{\mu} = 0$, pero el primal no es soluble.
\end{example}

\begin{example}[Óptimo primal existente, sin salto, pero dual no soluble]\label{ex:no_dual_sol}
    Sea el problema en $\Omega = \R$:
    \begin{mini*}{ }{x}{}{ }
        \addConstraint{x^2}{\le 0}
    \end{mini*}
    El único punto factible es $\bar{x} = 0$, por lo que $\bar{v} = f(0) = 0$. La Lagrangiana es $L(x, \mu) = x + \mu x^2$. Calculamos la función dual $\varphi(\mu) = \inf_{x \in \R} (x + \mu x^2)$:
    \begin{itemize}
        \item Si $\mu > 0$, el mínimo de la parábola se alcanza en $x = -1/(2\mu)$, dando $\varphi(\mu) = -1/(4\mu)$.
        \item Si $\mu = 0$, el ínfimo de la función lineal en $\R$ es $\varphi(0) = -\infty$.
    \end{itemize}
    El conjunto de factibilidad dual es el intervalo abierto $\Delta = (0, +\infty)$. El problema dual es $\max_{\mu > 0} (-1/(4\mu))$. Al aproximar $\mu \to \infty$, el supremo es $\bar{w} = 0$. Sin embargo, este valor no se alcanza para ningún multiplicador finito en el dominio abierto $\Delta$. El salto de dualidad es nulo, pero el problema dual no posee solución.
\end{example}

\begin{example}[Salto de dualidad estrictamente positivo]\label{ex:duality_gap}
    Tomemos el problema de optimización discreta sobre el conjunto no convexo $\Omega = \{0, 1\}$:
    \begin{mini*} {}{x_1 + x_2}{}{}
        \addConstraint{x_1^2 + x_2^2 - 2}{\le 0.}
    \end{mini*} 
    El único punto factible es $\bar{x} = 0$, lo que determina un valor óptimo primal $\bar{v} = f(0) = 0$.
    
    Para $\mu \ge 0$, la Lagrangiana es $L(x, \mu) = -x + \mu(x - 1/2)$. Evaluamos la función dual $\varphi(\mu)$ seleccionando el mínimo en el conjunto discreto $\Omega$:
    \[
        \varphi(\mu) = \min_{x \in \{0, 1\}} L(x, \mu) = \min \left\{ -\frac{1}{2}\mu, \, \frac{1}{2}\mu - 1 \right\}.
    \]
    El máximo de esta función cóncava y continua por partes se localiza en la intersección de ambas rectas:
    \[
        -\frac{1}{2}\mu = \frac{1}{2}\mu - 1 \implies \bar{\mu} = 1.
    \]
    El valor óptimo dual es $\bar{w} = \varphi(1) = -1/2$. Obtenemos un salto de dualidad estrictamente positivo:
    \[
        \bar{v} - \bar{w} = 0 - \left(-\frac{1}{2}\right) = \frac{1}{2} > 0.
    \]
\end{example}

% ---------------------------------------------------------
\subsection{Teoría de Multiplicadores de Lagrange Globales}
% ---------------------------------------------------------

\begin{definition}[Multiplicador de Lagrange Global]\label{def:lagrange_multiplier_global}
    Un vector $(\bar{\lambda}, \bar{\mu}) \in \R^l \times \R^m_+$ es un \textbf{multiplicador de Lagrange global} para el problema primal \eqref{eq:primal_general} si cumple:
    \begin{equation}
        \bar{v} = \inf_{x \in \Omega} L(x, \bar{\lambda}, \bar{\mu}).
        \label{eq:def_multiplier}
    \end{equation}
\end{definition}

Esta definición establece que las penalizaciones óptimas permiten hallar el valor óptimo del problema original mediante una minimización no restringida de la Lagrangiana sobre el conjunto base $\Omega$.

\begin{proposition}[Identidad entre Multiplicadores y Soluciones Duales]\label{prop:multipliers_dual_solutions}
    Supóngase que el valor óptimo primal $\bar{v}$ es finito y que el salto de dualidad es nulo. El conjunto de multiplicadores de Lagrange del problema primal coincide entonces con el conjunto de soluciones óptimas del problema dual.
\end{proposition}

\begin{proof}
    Si $(\bar{\lambda}, \bar{\mu})$ es un multiplicador de Lagrange, se tiene $\bar{\mu} \ge 0$ y $\varphi(\bar{\lambda}, \bar{\mu}) = \bar{v}$. Como el salto de dualidad es nulo, se cumple $\bar{v} = \bar{w}$, lo que da $\varphi(\bar{\lambda}, \bar{\mu}) = \bar{w}$. Esto demuestra que el par maximiza la función dual.
    
    Recíprocamente, si $(\bar{\lambda}, \bar{\mu}) \in \Delta$ es óptimo dual, entonces $\bar{\mu} \ge 0$ y $\varphi(\bar{\lambda}, \bar{\mu}) = \bar{w}$. Al cumplirse $\bar{w} = \bar{v}$, resulta inmediato que $\varphi(\bar{\lambda}, \bar{\mu}) = \bar{v}$, satisfaciendo la definición de multiplicador.
\end{proof}

\begin{example}[Identidad en un problema cuadrático bidimensional]\label{ex:identidad_multiplicadores}
    Consideremos la minimización de una función afín sobre un disco euclidiano compacto en $\Omega = \R^2$:
    \begin{mini*} {}{x_1 - x_2}{}{}
        \addConstraint{x_1 + x_2}{\le 0.}
    \end{mini*} 
    El mínimo de esta función sobre el disco de radio $\sqrt{2}$ se alcanza en $\bar{x} = (-1, -1)^\top$. En consecuencia, el \textbf{valor óptimo primal} está dado por $\bar{v} = f(-1, -1) = -2$.
    
    Para cualquier escalar $\mu \ge 0$, la Lagrangiana es $L(x, \mu) = x_1 + x_2 + \mu(x_1^2 + x_2^2 - 2)$. Si $\mu = 0$, $\varphi(0) = -\infty$. Para $\mu > 0$, la función es estrictamente convexa. Anulando el gradiente respecto a $x$:
    \[
        \nabla_x L(x, \mu) = \begin{pmatrix} 1 + 2\mu x_1 \\ 1 + 2\mu x_2 \end{pmatrix} = \begin{pmatrix} 0 \\ 0 \end{pmatrix}
        \implies x_1^*(\mu) = x_2^*(\mu) = -\frac{1}{2\mu}.
    \]
    Sustituyendo, obtenemos la función dual $\varphi(\mu) = -\frac{1}{2\mu} - 2\mu$.
    
    El problema dual exige maximizar esta función cóncava sobre el dominio $\mu > 0$. Derivando e igualando a cero: $\varphi'(\mu) = \frac{1}{2\mu^2} - 2 = 0 \implies \mu^2 = 1/4$. El único candidato admisible es $\bar{\mu} = 1/2$. Al evaluar, confirmamos que el \textbf{valor óptimo dual} es: $\bar{w} = \varphi(1/2) = -2$.
    
    Observamos empíricamente que $\bar{v} = \bar{w} = -2$ (brecha nula). El óptimo dual $\bar{\mu} = 1/2$ debe satisfacer la definición topológica de un multiplicador global (\cref{def:lagrange_multiplier_global}). Comprobémoslo:
    \[
        \inf_{x \in \R^2} L\left(x, \frac{1}{2}\right) = \inf_{x \in \R^2} \left[ \frac{1}{2}(x_1 + 1)^2 + \frac{1}{2}(x_2 + 1)^2 - 2 \right].
    \]
    Es evidente que el valor mínimo de esta suma de cuadrados es $-2$, alcanzado en $x = (-1,-1)^\top$. Hemos demostrado que el conjunto de multiplicadores de Lagrange y el de soluciones duales son idénticos.
\end{example}

\begin{theorem}[Garantía de Existencia de Solución Dual]\label{thm:dual_solvability_multiplier}
    Supóngase que el conjunto factible del problema primal es no vacío ($D \neq \emptyset$) y que el problema primal posee al menos un multiplicador de Lagrange global.
    \begin{enuthm}
        \item Si el conjunto de factibilidad dual es no vacío ($\Delta \neq \emptyset$), entonces el problema dual posee una solución óptima.
        \item Si el conjunto de factibilidad dual es vacío ($\Delta = \emptyset$), entonces el valor óptimo del problema primal es infinito negativo ($\bar{v} = -\infty$).
    \end{enuthm}
\end{theorem}

\begin{proof}
    Sea $(\bar{\lambda}, \bar{\mu}) \in \R^l \times \R^m_+$ un multiplicador de Lagrange. Por definición (\cref{def:lagrange_multiplier_global}), $\varphi(\bar{\lambda}, \bar{\mu}) = \bar{v}$.
    
    Si $\Delta \neq \emptyset$, por el Teorema de Dualidad Débil (\cref{thm:dualidad_debil_general}), $\bar{v}$ acota superiormente a cualquier valor factible dual. Puesto que el par $(\bar{\lambda}, \bar{\mu})$ alcanza este límite superior, pertenece a $\Delta$ y es un maximizador global de la función dual.
    
    Si $\Delta = \emptyset$, implica que para todo par $(\lambda, \mu)$ se tiene $\varphi(\lambda, \mu) = -\infty$. Evaluando en el multiplicador garantizado por hipótesis, obtenemos $\varphi(\bar{\lambda}, \bar{\mu}) = -\infty$. Por ende, $\bar{v} = -\infty$.
\end{proof}

\begin{example}[Infactibilidad dual y divergencia primal]\label{ex:teorema_existencia_dual}
    Consideremos el siguiente problema lineal sobre $\Omega = \R^2$:
    \begin{mini*} {}{x_1 - x_2}{}{}
        \addConstraint{x_1 + x_2}{\le 0.}
    \end{mini*}
    El dominio factible $D$ contiene a $(0,0)$. Tomando $x(t) = (-t, t)^\top$ para $t > 0$, la función objetivo evaluada en esta trayectoria es $-2t$. Al hacer $t \to \infty$, el problema es ilimitado: $\bar{v} = -\infty$.
    
    Construyamos la función Lagrangiana para $\mu \ge 0$: $L(x, \mu) = x_1(1 + \mu) + x_2(\mu - 1)$. Para que el ínfimo sea finito, se requeriría $1 + \mu = 0$ y $\mu - 1 = 0$, lo cual es una contradicción. Concluimos que para cualquier $\mu \ge 0$, $\varphi(\mu) = -\infty$. Por tanto, $\Delta = \emptyset$, lo que corrobora el Teorema \ref{thm:dual_solvability_multiplier}.
\end{example}

\begin{theorem}[Condiciones de Kuhn-Tucker No Diferenciables]\label{thm:kuhn_tucker_non_diff}
    Un punto $\bar{x} \in D$ es solución óptima del problema primal y $(\bar{\lambda}, \bar{\mu})$ es un multiplicador de Lagrange si y solo si se cumplen las condiciones:
    \begin{align}
        L(\bar{x}, \bar{\lambda}, \bar{\mu}) & = \min_{x \in \Omega} L(x, \bar{\lambda}, \bar{\mu}), \label{eq:kt_min_global} \\
        \bar{\mu}_i g_i(\bar{x})             & = 0, \quad \forall i = 1, \dots, m. \label{eq:kt_complementarity_non_diff}
    \end{align}
\end{theorem}

\begin{proof}
    \( (\implies) \) Como $\bar{x}$ es óptimo y $(\bar{\lambda}, \bar{\mu})$ es multiplicador, se tiene $f(\bar{x}) = \inf_{x \in \Omega} L(x, \bar{\lambda}, \bar{\mu})$. Como $\bar{x} \in D$, resulta $L(\bar{x}, \bar{\lambda}, \bar{\mu}) = f(\bar{x}) + \interno{\bar{\mu}}{g(\bar{x})}$. La condición $\bar{\mu} \ge 0, g(\bar{x}) \le 0$ exige que $\interno{\bar{\mu}}{g(\bar{x})} \le 0$. Al ser el ínfimo $f(\bar{x})$, se debe verificar de forma obligatoria que $L(\bar{x}, \bar{\lambda}, \bar{\mu}) = f(\bar{x})$, lo que prueba \eqref{eq:kt_min_global} y requiere $\bar{\mu}_i g_i(\bar{x}) = 0$.
    
    \( (\impliedby) \) Por holgura complementaria y factibilidad, se tiene $L(\bar{x}, \bar{\lambda}, \bar{\mu}) = f(\bar{x})$. Sustituyendo en la minimización global \eqref{eq:kt_min_global}, resulta $f(\bar{x}) = \inf_{x \in \Omega} L(x, \bar{\lambda}, \bar{\mu})$, lo cual demuestra que $(\bar{\lambda}, \bar{\mu})$ es un multiplicador y $\bar{x}$ es óptimo global.
\end{proof}

% ---------------------------------------------------------
\subsection{Dualidad Fuerte y la Condición de Slater}
% ---------------------------------------------------------

\begin{definition}[Condición de Regularidad de Slater]\label{def:slater}
    Se dice que el problema primal \eqref{eq:primal_general} satisface la \textbf{condición de Slater} si las restricciones de igualdad son afines ($h(x) = Ax - a$), las restricciones de desigualdad $g_i$ son funciones convexas en $\Omega$, y existe un punto de regularidad $\hat{x} \in \interior \Omega$ tal que:
    \begin{equation}
        A\hat{x} - a = 0 \quad \text{y} \quad g_i(\hat{x}) < 0 \quad \forall i = 1, \dots, m.
    \end{equation}
\end{definition}

\begin{theorem}[Teorema de Dualidad Fuerte de Slater]\label{thm:dualidad_fuerte_slater}
    Supóngase que en el problema primal \eqref{eq:primal_general} las funciones $f$ y $g_i$ son convexas, las restricciones de igualdad son afines dadas por $h(x) = Ax - a$ con $A \in \R^{l \times n}$, y el valor óptimo primal $\bar{v}$ es finito. Si se satisface la condición de regularidad de Slater (\cref{def:slater}), entonces el salto de dualidad es nulo ($\bar{v} = \bar{w}$), y el problema dual admite al menos una solución óptima $(\bar{\lambda}, \bar{\mu}) \in \Delta$.
\end{theorem}

\begin{proof}
    Supongamos, sin pérdida de generalidad, que las filas de la matriz $A$ son linealmente independientes. Definamos el conjunto de valores admisibles para la función objetivo y las perturbaciones de las restricciones primales:
    \[
        C = \left\{ (y, z, c) \in \R \times \R^m \times \R^l \mathrel{\bigg|} \exists x \in \Omega \text{ t.q. } f(x) \le y, \, g(x) \le z, \, Ax - a = c \right\}.
    \]
    La convexidad de $f$ y $g_i$, sumada al carácter afín de $h$, garantizan que $C$ es un conjunto convexo en $\R \times \R^m \times \R^l$. Dado que $\bar{v}$ es el valor óptimo primal, el punto $(\bar{v}, 0, 0)$ se encuentra en la frontera del conjunto convexo $C$.
    
    Aplicamos el Teorema de Separación de Conjuntos Convexos (Lema de Minkowski). Existe un hiperplano de soporte en $(\bar{v}, 0, 0)$, caracterizado por un vector normal no nulo $(\beta, \bar{\mu}, \bar{\lambda}) \in \R \times \R^m \times \R^l$, tal que:
    \begin{equation}
        \beta \bar{v} \le \beta y + \interno{\bar{\mu}}{z} + \interno{\bar{\lambda}}{c} \quad \forall (y, z, c) \in C.
        \label{eq:slater_separacion}
    \end{equation}
    Para que la desigualdad se cumpla asintóticamente es necesario que $\beta \ge 0$ y $\bar{\mu} \ge 0$.
    
    Asumamos que $\beta = 0$. Evaluando \eqref{eq:slater_separacion} en el punto factible de Slater $\hat{x}$, se obtiene $0 \le \interno{\bar{\mu}}{g(\hat{x})}$. Dado que $\bar{\mu} \ge 0$ y cada componente $g_i(\hat{x}) < 0$, la única posibilidad es que $\bar{\mu} = 0$. Sustituyendo $\beta = 0$ y $\bar{\mu} = 0$, nos queda $0 \le \interno{\bar{\lambda}}{c}$ para todo $c$ factible. Como $A$ tiene rango completo, esto es posible solo si $\bar{\lambda} = 0$. Esto implica que el vector normal es nulo, contradiciendo el Teorema de Separación. Por consiguiente, $\beta > 0$.
    
    Dividiendo \eqref{eq:slater_separacion} por $\beta$ y definiendo $\lambda^* = \bar{\lambda}/\beta \in \R^l$ y $\mu^* = \bar{\mu}/\beta \in \R^m_+$. Tomando $(y, z, c) = (f(x), g(x), Ax - a) \in C$, la inecuación se reduce a:
    \[
        \bar{v} \le f(x) + \interno{\mu^*}{g(x)} + \interno{\lambda^*}{Ax - a} = L(x, \lambda^*, \mu^*).
    \]
    Tomando el ínfimo sobre todo $x \in \Omega$, concluimos que $\bar{v} \le \varphi(\lambda^*, \mu^*)$. A su vez, por Dualidad Débil (\cref{thm:dualidad_debil_general}), sabemos que $\varphi(\lambda^*, \mu^*) \le \bar{v}$. Así, $\varphi(\lambda^*, \mu^*) = \bar{v}$, probando que no hay salto de dualidad y el óptimo dual se alcanza.
\end{proof}

\begin{example}[Verificación de la dualidad fuerte]\label{ex:dualidad_fuerte_slater}
    Sea el problema de optimización en $\Omega = \R$:
    \begin{mini*} {}{(x - 3)^2}{}{}
        \addConstraint{x - 2}{\leq 0}
    \end{mini*}
    El origen $\hat{x} = 0 \in \interior \R$ cumple $g(0) = -2 < 0$, satisfaciendo la condición de Slater. El mínimo restringido se alcanza en la frontera $\bar{x} = 2$, dando $\bar{v} = 1$.
    
    Construimos la función dual para $\mu \ge 0$: $\varphi(\mu) = \inf_{x \in \R} \{ (x - 3)^2 + \mu(x - 2) \} = -\frac{\mu^2}{4} + \mu$. Maximizando la función dual, calculamos el óptimo $\bar{\mu} = 2$. El valor óptimo dual es $\bar{w} = \varphi(2) = 1$. Verificamos que $\bar{v} = \bar{w} = 1$.
\end{example}

% ---------------------------------------------------------
\subsection{El Problema Dual de Wolfe}
% ---------------------------------------------------------

\begin{definition}[Problema Dual de Wolfe]\label{def:wolfe_dual}
    Asociado al problema primal \eqref{eq:primal_general}, asumiendo que $\Omega$ es abierto en $\Rn$ y que $f, g, h$ son continuamente diferenciables, el \textbf{problema dual de Wolfe} se define como:
    \begin{equation}
        \max_{\substack{x \in \Omega, \lambda \in \R^l \\ \mu \ge 0}} L(x, \lambda, \mu) \quad \text{sujeto a} \quad \nabla_x L(x, \lambda, \mu) = 0.
        \label{eq:wolfe_dual}
    \end{equation}
    Denotaremos al conjunto de factibilidad de este problema como $W$.
\end{definition}

\begin{theorem}[Dualidad Débil de Wolfe]\label{thm:dualidad_debil_wolfe}
    Si el problema primal original es de minimización convexa (es decir, $f$ y $g_i$ son convexas, y $h$ es afín), entonces para cualquier punto factible primal $\tilde{x} \in D$ y cualquier punto factible del dual de Wolfe $(x, \lambda, \mu) \in W$, se verifica que:
    \begin{equation}
        L(x, \lambda, \mu) \le f(\tilde{x}).
    \end{equation}
\end{theorem}

\begin{proof}
    Dado que $f$ y $g_i$ son convexas y $h$ es afín, la Lagrangiana $x \mapsto L(x, \lambda, \mu)$ preserva la convexidad. Por la definición de $W$, sabemos que $\nabla_x L(x, \lambda, \mu) = 0$. En funciones convexas, la anulación del gradiente garantiza que $x$ es un minimizador global: $L(x, \lambda, \mu) = \inf_{z \in \Omega} L(z, \lambda, \mu) = \varphi(\lambda, \mu)$. Aplicando el Teorema de Dualidad Débil General (\cref{thm:dualidad_debil_general}), sabemos que $\varphi(\lambda, \mu) \le f(\tilde{x})$. Sustituyendo la igualdad previa, se concluye: $L(x, \lambda, \mu) \le f(\tilde{x})$.
\end{proof}

\begin{example}[Coincidencia del Dual de Wolfe en Programación Lineal]\label{ex:wolfe_lp}
    Sea el programa primal lineal: $\min \interno{c}{x}$ sujeto a $Bx \ge b$. Adaptando a $g(x) = b - Bx \le 0$, la Lagrangiana es $L(x, \mu) = \interno{c - B^\top \mu}{x} + \interno{b}{\mu}$. La restricción del dual de Wolfe es $\nabla_x L(x, \mu) = 0$, lo que equivale a $c - B^\top \mu = 0$. Evaluando la función objetivo $L(x, \mu)$ bajo esta condición, el término con $x$ se anula, resultando $L(x, \mu) = \interno{b}{\mu}$. El dual de Wolfe se simplifica a maximizar $\interno{b}{\mu}$ sujeto a $B^\top \mu = c$ y $\mu \ge 0$, que es la formulación dual estándar planteada en \eqref{eq:dual_lp}.
\end{example}

% ---------------------------------------------------------
\subsection{Caracterización mediante Puntos de Silla}
% ---------------------------------------------------------

\begin{definition}[Punto de Silla]\label{def:punto_silla}
    Un punto $(\bar{x}, \bar{\lambda}, \bar{\mu}) \in \Omega \times \R^l \times \R^m_+$ es un \textbf{punto de silla} de la Lagrangiana si verifica la doble desigualdad topológica:
    \begin{equation}
        L(\bar{x}, \lambda, \mu) \le L(\bar{x}, \bar{\lambda}, \bar{\mu}) \le L(x, \bar{\lambda}, \bar{\mu})
        \label{eq:punto_silla}
    \end{equation}
    para todo $x \in \Omega$, todo $\lambda \in \R^l$ y todo $\mu \ge 0$.
\end{definition}

\begin{theorem}[Teorema del Punto de Silla]\label{thm:teorema_punto_silla}
    Un punto $\bar{x} \in \Omega$ es solución óptima del problema primal y $(\bar{\lambda}, \bar{\mu})$ es solución óptima del problema dual asociado con salto de dualidad nulo si y solo si $(\bar{x}, \bar{\lambda}, \bar{\mu})$ es un punto de silla de la Lagrangiana.
\end{theorem}

\begin{proof}
    \( (\implies ): \) Por el \cref{lem:sup_lagrangiana} y la factibilidad de $\bar{x}$, se tiene $L(\bar{x}, \lambda, \mu) \le f(\bar{x}) = L(\bar{x}, \bar{\lambda}, \bar{\mu})$, probando la primera desigualdad. Para la segunda, la dualidad fuerte implica $L(\bar{x}, \bar{\lambda}, \bar{\mu}) = \varphi(\bar{\lambda}, \bar{\mu}) \le L(x, \bar{\lambda}, \bar{\mu})$.
    
    \( (\impliedby ): \) Si es un punto de silla, al tomar el supremo en las variables duales de la primera desigualdad: $\sup L(\bar{x}, \lambda, \mu) \le L(\bar{x}, \bar{\lambda}, \bar{\mu}) < +\infty$. Por el \cref{lem:sup_lagrangiana}, esto fuerza a que $\bar{x} \in D$ y $L(\bar{x}, \bar{\lambda}, \bar{\mu}) = f(\bar{x})$. De la segunda desigualdad, $f(\bar{x}) \le L(x, \bar{\lambda}, \bar{\mu})$. Tomando el ínfimo sobre $x \in \Omega$, resulta $f(\bar{x}) \le \varphi(\bar{\lambda}, \bar{\mu})$. Por el Teorema de Dualidad Débil, esto obliga a que $f(\bar{x}) = \varphi(\bar{\lambda}, \bar{\mu})$, anulando el salto de dualidad.
\end{proof}

\begin{example}[Punto de silla de la Lagrangiana]\label{ex:punto_silla_verificacion}
    Tomemos el problema: $\min x^2 - 4x$ sujeto a $x - 1 \le 0$. El mínimo restringido en $D = (-\infty, 1]$ se alcanza en $\bar{x} = 1$, con $\bar{v} = -3$. La Lagrangiana es $L(x, \mu) = x^2 - 4x + \mu(x - 1)$. Igualando a cero la derivada en $\bar{x} = 1$, obtenemos $\bar{\mu} = 2$.
    
    Verificamos las condiciones de punto de silla para $(\bar{x}, \bar{\mu}) = (1, 2)$:
    \begin{enuthm}
        \item \textbf{Izquierda:} Para cualquier $\mu \ge 0$, $L(1, \mu) = -3 \le L(1, 2) = -3$, lo cual se cumple.
        \item \textbf{Derecha:} Para cualquier $x \in \R$, $L(1, 2) \le L(x, 2) \implies -3 \le x^2 - 2x - 2 \implies (x - 1)^2 \ge 0$, la cual se verifica siempre.
    \end{enuthm}
    Esto confirma que la terna constituye un punto de silla de la Lagrangiana.
\end{example}

% ---------------------------------------------------------
\subsection{Subgradiente de la Función Dual}
% ---------------------------------------------------------

\begin{proposition}[Subgradiente de la Función Dual]\label{prop:dual_subgradient}
    Sean $( \lambda, \mu) \in \Delta$ y sea $x( \lambda, \mu) \in \Omega$ tal que se alcanza el ínfimo parcial, es decir, $\varphi( \lambda, \mu) = L(x( \lambda, \mu), \lambda, \mu)$. El vector de violaciones funcionales:
    \begin{equation}
        -\begin{pmatrix} h(x( \lambda, \mu)) \\ g(x( \lambda, \mu)) \end{pmatrix} \in \partial (-\varphi)( \lambda, \mu)
    \end{equation}
    es un subgradiente de la función dual negativa $-\varphi$ en el punto $(\lambda, \mu)$.
\end{proposition}

\begin{proof}
    Un vector $s$ es un subgradiente si $\varphi(\eta, \nu) \le \varphi(\lambda, \mu) - \interno{s}{( \eta - \lambda , \nu - \mu )^\top}$. Partiendo de la definición de la función dual en un punto genérico $(\eta, \nu) \in \Delta$:
    \[
        \varphi( \eta, \nu) = \inf_{z \in \Omega} L(z, \eta, \nu) \le L(x( \lambda, \mu), \eta, \nu).
    \]
    Desarrollando la expresión de la Lagrangiana evaluada en $x( \lambda, \mu)$:
    \[
        L(x(\lambda, \mu), \eta, \nu) = \varphi(\lambda, \mu) + \interno{\eta - \lambda}{h(x( \lambda, \mu))} + \interno{\nu - \mu}{g(x( \lambda, \mu))}.
    \]
    Esto verifica la desigualdad del subgradiente al tomar $s = -(h(x( \lambda, \mu)), g(x( \lambda, \mu)))^\top$.
\end{proof}

\begin{controlbox}[Interpretación de puntos de silla]
    Si se encuentra un punto de silla para la función Lagrangiana de un problema no convexo, ¿se puede asegurar que dicho punto define la solución global del problema primal? Justifique su respuesta basándose en las hipótesis del \cref{thm:teorema_punto_silla}.
\end{controlbox}

\begin{notabox}[Relajación Lagrangiana]
    El Teorema de Dualidad Débil sustenta la técnica conocida como \textbf{relajación Lagrangiana}, utilizada en optimización entera y combinatoria. Al trasladar restricciones difíciles a la función objetivo mediante multiplicadores fijos $\mu \ge 0$, se obtiene un problema relajado fácil de resolver cuyo óptimo proporciona siempre una cota inferior rigurosa para el valor mínimo del problema original.
\end{notabox}